% !TEX program = pdflatex
% Hugging Face Deep RL Course - Bonus Unit 2 (complete MDX study edition).
% Compile with latexmk -pdf, from notes/en/, using the unchanged ../preamble.tex.
\documentclass[11pt,a4paper]{article}

% Shared typography/layout for Hugging Face Deep RL Course notes.
% Compile with pdfLaTeX / latexmk -pdf with notes/ as the working directory.
% e.g. cd notes && latexmk -pdf unit1.tex

\usepackage[T1]{fontenc}
\usepackage[utf8]{inputenc}
\usepackage[english]{babel}
\usepackage{lmodern}
\usepackage{textcomp}
\usepackage{geometry}
\geometry{a4paper,margin=22mm,headheight=14pt,headsep=8mm,footskip=12mm}

\usepackage{amsmath,amssymb,mathtools}
\usepackage{graphicx,caption,needspace,placeins}
\usepackage{array,booktabs,tabularx}
\usepackage{xcolor}
\usepackage{microtype}
\usepackage{enumitem}
\usepackage{listings}
\usepackage[most]{tcolorbox}
\usepackage{titlesec}
\usepackage{fancyhdr}
\usepackage{hyperref}
\usepackage{bookmark}

\definecolor{ink}{HTML}{243041}
\definecolor{accent}{HTML}{38598C}
\definecolor{secondary}{HTML}{60768A}
\definecolor{codebg}{HTML}{F4F6F8}
\definecolor{codeborder}{HTML}{E3E8EE}
\definecolor{notebg}{HTML}{F3F7FC}

\hypersetup{
  colorlinks=true,
  linkcolor=accent,
  urlcolor=accent,
  citecolor=accent,
  pdfborder={0 0 0}
}
\urlstyle{same}

\setlength{\parindent}{0pt}
\setlength{\parskip}{0.38em plus 0.12em minus 0.05em}

\setlist[itemize]{
  topsep=0.35em,
  itemsep=0.13em,
  leftmargin=1.7em
}
\setlist[enumerate]{
  topsep=0.35em,
  itemsep=0.13em,
  leftmargin=1.8em
}

\setcounter{tocdepth}{2}
\setcounter{secnumdepth}{2}

\titleformat{\section}
  {\Large\bfseries\color{ink}}
  {\thesection}{0.7em}{}

\titleformat{\subsection}
  {\large\bfseries\color{accent}}
  {\thesubsection}{0.7em}{}

\titleformat{\subsubsection}
  {\normalsize\bfseries\color{ink}}
  {\thesubsubsection}{0.6em}{}

\titlespacing*{\section}
  {0pt}{2.2ex plus 0.7ex}{0.9ex}
\titlespacing*{\subsection}
  {0pt}{1.7ex plus 0.4ex}{0.6ex}
\titlespacing*{\subsubsection}
  {0pt}{1.3ex plus 0.3ex}{0.45ex}

\pagestyle{fancy}
\fancyhf{}

\fancyhead[L]{
  \footnotesize\sffamily\color{secondary}
  HF Deep RL Course \enspace/\enspace Unit 1
}
\fancyhead[R]{
  \footnotesize\sffamily\color{secondary}
  \nouppercase{\rightmark}
}
\fancyfoot[C]{
  \footnotesize\color{secondary}\thepage
}
\renewcommand{\headrulewidth}{0.25pt}

\captionsetup{
  font=small,
  labelfont={bf,color=secondary},
  textfont={color=ink},
  justification=centering,
  skip=4pt,
  hypcap=false
}

\lstdefinestyle{hfcode}{
  basicstyle=\ttfamily\footnotesize,
  backgroundcolor=\color{codebg},
  frame=single,
  rulecolor=\color{codeborder},
  numbers=none,
  showstringspaces=false,
  breaklines=true,
  breakatwhitespace=false,
  columns=fullflexible,
  keepspaces=true,
  xleftmargin=0pt,
  xrightmargin=0pt,
  aboveskip=0.6em,
  belowskip=0.6em,
  framexleftmargin=7pt,
  framexrightmargin=7pt,
  framextopmargin=5pt,
  framexbottommargin=5pt,
  upquote=true,
  tabsize=4,
}
\lstset{style=hfcode}

\newcommand{\CourseFigure}[2]{%
  \par\addvspace{0.55em}\noindent%
  \begin{minipage}{\linewidth}%
    \centering
    \includegraphics[
      width=0.86\linewidth,
      height=0.295\textheight,
      keepaspectratio
    ]{#1}\par
    \captionof{figure}{#2}%
  \end{minipage}%
  \par\addvspace{0.55em}%
}

\newtcolorbox{coursenote}{
  enhanced,
  breakable,
  colback=notebg,
  colframe=accent!30,
  boxrule=0.35pt,
  arc=1.2mm,
  left=9pt,
  right=9pt,
  top=6pt,
  bottom=6pt
}

\renewenvironment{quote}
  {\begin{coursenote}}
  {\end{coursenote}}

\emergencystretch=2.5em

\providecommand{\tightlist}{
  \setlength{\itemsep}{0pt}
  \setlength{\parskip}{0pt}
}

\providecommand{\passthrough}[1]{#1}

\Urlmuskip=0mu plus 6mu\relax

\graphicspath{{../assets/unitbonus2/}}
\fancyhead[L]{\footnotesize\sffamily\color{secondary}HF Deep RL Course \enspace/\enspace Bonus Unit 2}
\fancyhead[R]{\footnotesize\sffamily\color{secondary}Hyperparameter Tuning}
\title{\textbf{Deep Reinforcement Learning}\\[0.35em]\large Bonus Unit 2: Automatic Hyperparameter Tuning with Optuna}
\author{Hugging Face Deep RL Course\\\small Complete MDX-to-\LaTeX{} study edition}
\date{}

\begin{document}
\maketitle
\vspace{-1em}
\begin{center}\small\color{secondary}Source: Hugging Face Deep RL Course (original English text).\end{center}
\vspace{-0.45em}
\begingroup\small\setlength{\parskip}{0pt}\tableofcontents\endgroup
\vspace{0.3em}

\section{Introduction}\label{introduction}

One of the most critical tasks in Deep Reinforcement Learning is to \textbf{find a good set of training hyperparameters}.

\CourseFigure{optuna-logo.png}{Optuna library logo. \protect\href{https://raw.githubusercontent.com/optuna/optuna/master/docs/image/optuna-logo.png}{[original image]}}

\href{https://optuna.org/}{Optuna} is a library that helps you to automate the search. In this Unit, we'll study a \textbf{little bit of the theory behind automatic hyperparameter tuning}. We'll first try to optimize the parameters of the DQN studied in the last unit manually. We'll then \textbf{learn how to automate the search using Optuna}.

\section{Optuna Tutorial}\label{optuna}

The content below comes from \href{https://araffin.github.io/tools-for-robotic-rl-icra2022/}{Antonin's Raffin ICRA 2022 presentations}, he's one of the founders of Stable-Baselines and RL-Baselines3-Zoo.

\subsection{The theory behind Hyperparameter tuning}

\begin{coursenote}
  \textbf{Video lecture: The theory behind Hyperparameter tuning}\par
  \href{https://www.youtube.com/watch?v=AidFTOdGNFQ}{Watch on YouTube (AidFTOdGNFQ)}.
\end{coursenote}

\subsection{Optuna Tutorial}

\begin{coursenote}
  \textbf{Video tutorial: Optuna}\par
  \href{https://www.youtube.com/watch?v=ihP7E76KGOI}{Watch on YouTube (ihP7E76KGOI)}.
\end{coursenote}

The notebook $\rightarrow$ \href{https://colab.research.google.com/github/araffin/tools-for-robotic-rl-icra2022/blob/main/notebooks/optuna_lab.ipynb}{here}.

\section{Hands-on}\label{hands-on}

Now that you've learned to use Optuna, here are some ideas to apply what you've learned:

\begin{enumerate}[label=\textbf{\arabic*.},leftmargin=1.8em,itemsep=0.6em]
  \item \textbf{Beat your LunarLander-v2 agent results}, by using Optuna to find a better set of hyperparameters. You can also try with another environment, such as MountainCar-v0 and CartPole-v1.
  \item \textbf{Beat your SpaceInvaders agent results}.
\end{enumerate}

By doing this, you'll see how valuable and powerful Optuna can be in training better agents.

Have fun!

Finally, we would love \textbf{to hear what you think of the course and how we can improve it}. If you have some feedback then please $\rightarrow$ \href{https://forms.gle/BzKXWzLAGZESGNaE9}{fill out this form}.

\subsection*{Keep Learning, stay awesome}

\end{document}
